\documentclass[12pt]{article}
\usepackage[T1]{fontenc}
\usepackage[utf8]{inputenc}
\usepackage{lmodern}
\usepackage{textcomp}
\usepackage{enumitem}
\usepackage{geometry}
\usepackage{graphicx}
\usepackage{amsmath, amssymb}
\geometry{margin=1in}
\begin{document}
\begin{center}
History - Undergraduate Level Assessment 16 \
Total Marks: 40 \
Answer all questions.
\end{center}

\section*{Multiple Choice (10 marks)}
\begin{enumerate}[label=\arabic*.]
\item Inca architecture is noted for its:
\begin{enumerate}[label=(\alph*)]
    \item focus on creating intricate geometric patterns.
    \item construction of structures exclusively on hilltops.
    \item employment of advanced hydraulic systems.
    \item incorporation of wood, earth, and stone.
    \item reliance on natural materials such as leaves and vines.
\end{enumerate}
\item The \_\_\_\_\_\_\_\_\_\_ tradition features finely made, leaf-shaped stone blades. The \_\_\_\_\_\_\_\_\_\_ tradition, which followed behind it, emphasizes bone and antler work.
\begin{enumerate}[label=(\alph*)]
    \item Aurignacian; Solutrean
    \item Gravettian; Solutrean
    \item Gravettian; Aurignacian
    \item Gravettian; Magdelanian
    \item Magdelanian; Gravettian
\end{enumerate}
\item Analysis of the teeth of a hominid species discovered in Malapa Cave in South Africa, dating just after 2 million years ago, indicates a diet that consisted primarily of:
\begin{enumerate}[label=(\alph*)]
    \item meat, fish, and some plants.
    \item primarily C 3 plants that included barley and oats.
    \item large quantities of insects and small mammals.
    \item a balanced diet of fruits, vegetables, meat, and fish.
    \item primarily C 4 plants that included corn and potatoes.
\end{enumerate}
\item In primates, the location of which of the following determines if it is a quadruped or biped?
\begin{enumerate}[label=(\alph*)]
    \item the foramen magnum
    \item the skull
    \item the feet
    \item the femur
    \item the spine
\end{enumerate}
\item At its peak, the palace at Knossos is thought to have had over \_\_\_\_\_\_\_\_\_ rooms.
\begin{enumerate}[label=(\alph*)]
    \item 100
    \item 500
    \item 1,000
    \item 5,000
\end{enumerate}
\end{enumerate}

\section*{True / False (10 marks)}
\textit{Answer True or False}
\begin{enumerate}[label=\arabic*.]
\setcounter{enumi}{5}
\item True or False: In China, urbanization increased as the population grew and as the division of labor grew more complex. Large urban centers, such as Nanjing and Beijing, also contributed to the a different answer.
\item True or False: She made her theatrical debut in a production of David Rabe's Goose and Tom-Tom; the film and play both co-starred Penn. The next year, Madonna was featured in the film Who's That Girl. She contributed four songs to its soundtrack, including the title track and "a different answer".
\item True or False: In a 'Letter to American People' written by Osama bin Laden in 2002, he stated that one of the reasons he was fighting America is because of its support of India on the Kashmir issue. While on a trip to Delhi in 2002, U.S. Secretary of Defense Donald Rumsfeld suggested that Al-Qaeda was active in Kashmir, though he did not have any hard evidence.
\item True or False: The names Achaioi and Danaoi seem to be pre-Dorian belonging to the people who were overthrown. They were forced to the region that later bore the name Achaea after the Dorian invasion. In the 5 th century BC, they were redefined as contemporary speakers of Aeolic Greek which was spoken mainly in Thessaly, Boeotia and Lesbos.
\item True or False: From the early 6 th century they spread to inhabit most of Central and Eastern Europe and Southeast Europe, whilst Slavic mercenaries fighting for the Byzantines and Arabs settled Asia Minor and even as far as Syria. a different answer[better source needed] Presently over half of Europe's territory is inhabited by Slavic-speaking communities, but every Slavic ethnicity has emigrated to other continents.
\end{enumerate}

\section*{Long Form Response (20 marks)}
\textit{Answer each question with a clear and complete explanation.}
\begin{enumerate}[label=\arabic*.]
\setcounter{enumi}{10}
\item Read the following passage and answer the question: Bacteria were first observed by the Dutch microscopist Antonie van Leeuwenhoek in 1676, using a single-lens microscope of his own design. He then published his observations in a series of letters to the Royal Society of London. Bacteria were Leeuwenhoek's most remarkable microscopic discovery. They were just at the limit of what his simple lenses could make out and, in one of the most striking hiatuses in the history of science, no one else would see them again for over a century. Only then were his by-then-largely-forgotten observations of bacteria - as opposed to his famous "animalcules" (spermatozoa) - taken seriously. Question: When was the first time bacteria were observed?
\item Read the following passage and answer the question: In 1358, the Sakya viceregal regime installed by the Mongols in Tibet was overthrown in a rebellion by the Phagmodru myriarch Tai Situ Changchub Gyaltsen (1302-1364). The Mongol Yuan court was forced to accept him as the new viceroy, and Changchub Gyaltsen and his successors, the Phagmodrupa Dynasty, gained de facto rule over Tibet. Question: Which dynasty became ruler of Tibet?
\end{enumerate}

\vfill
\noindent\textit{End of Paper}
\end{document}